\begin{frame}{What we measured}
    \begin{table}[ht!]
  \centering
  \begin{tabular}{lccccc}
    \hline
    config & $d$ & $N$ & boundary condition & perturbation & $n$ \\
    \hline
    A & 2 & 40 & open   & nonconservative & $10^7$ \\
    B & 2 & 40 & open   & conservative    & $10^7$ \\
    C & 2 & 40 & closed & nonconservative & $10^7$ \\
    D & 3 & 20 & open   & nonconservative & $10^7$ \\
    E & 3 & 20 & open   & conservative    & $10^7$ \\
    F & 3 & 20 & closed & nonconservative & $10^7$ \\
    G & 4 & 20 & open   & nonconservative & $10^6$ \\
    H & 4 & 20 & open   & conservative    & $10^6$ \\
    I & 4 & 20 & closed & nonconservative & $10^6$ \\
    \hline
  \end{tabular}
  \caption{Simulated model configurations and measurement counts used in this work.}
  \label{tab:model_configurations}
\end{table}
\end{frame}

\begin{frame}{Did we observe power laws ?}
    \pause
    \begin{figure}
        \centering
        \includegraphics[width=\linewidth,height=0.55\textheight,keepaspectratio]{../../figures/loglog_plots/N40d2_open_conservative_lambda}
        \caption{Measured probability distribution of $l$ for a sandpile of dimension $N=40,\,d=2$ with open boundaries and conservative perturbation.}
    \end{figure}
\end{frame}

\begin{frame}{Did we observe power laws ?}

    \begin{figure}
        \centering
        \includegraphics[width=\linewidth,height=0.55\textheight,keepaspectratio]{../../figures/loglog_plots/N40d2_open_conservative_tau}
        \caption{Measured probability distribution of $s$ for a sandpile of dimension $N=40,\,d=2$ with open boundaries and conservative perturbation.}
    \end{figure}
\end{frame}


\begin{frame}{Did we observe power laws ?}
 
    \begin{figure}
        \centering
        \includegraphics[width=\linewidth,height=0.55\textheight,keepaspectratio]{../../figures/loglog_plots/N40d2_open_conservative_alpha}
        \caption{Measured probability distribution of $t$ for a sandpile of dimension $N=40,\,d=2$ with open boundaries and conservative perturbation.}
    \end{figure}
\end{frame}



\begin{frame}{Resulting exponentes}
    \begin{table}[!ht]
    \centering
    \begin{tabular}{cccc}
\toprule
model & $\tau$ & $\alpha$ & $\lambda$ \\
\midrule
A & $ (2.0182\pm0.0077) $ & $ (2.058\pm0.035) $ & $ (1.61\pm0.15) $ \\
D & $ (2.287\pm0.023) $ & $ (2.466\pm0.087) $ & $ (2.27\pm0.27) $ \\
G & $ (2.397\pm0.032) $ & $ (2.62\pm0.13) $ & $ (2.83\pm0.51) $ \\
\bottomrule
\end{tabular}

    \caption{Scaling exponents of the probability density functions for the open-boundary, nonconservative sandpile configurations (A, D, G). The exponents $\alpha$, $\tau$, and $\lambda$ are obtained from weighted linear least-squares fits on double-logarithmic PDF plots in the selected scaling window; The errors were determined by the aforementioned procedure.}
    \label{tab:scaling_exp_open_nonconservative}
    \end{table}
\end{frame}

\begin{frame}{Comparing the exponentes to literature}
    \begin{table}[!ht]
    \centering
    \begin{tabular}{cccc}
\toprule
model & $\tau$ & $\alpha$ & $\lambda$ \\
\midrule
A & $ (2.0182\pm0.0077) $ & $ (2.058\pm0.035) $ & $ (1.61\pm0.15) $ \\
D & $ (2.287\pm0.023) $ & $ (2.466\pm0.087) $ & $ (2.27\pm0.27) $ \\
G & $ (2.397\pm0.032) $ & $ (2.62\pm0.13) $ & $ (2.83\pm0.51) $ \\
\bottomrule
\end{tabular}

    \end{table}
    \begin{figure}
        \centering
        \includegraphics[width=\linewidth,height=0.55\textheight,keepaspectratio]{values_from_paper}
        \caption*{\cite{christensen1991dynamical}}
    \end{figure}
\end{frame} 

